\section{Evaluasi}
\label{sec:evaluasi-bab13}

\subsection{Soal Pilihan Ganda}

\begin{enumerate}
    \item Instruksi \asm{CALL} \eng{near} menyimpan \dots ke dalam
          stack.
          \begin{enumerate}
              \item register \reg{CS} dan \reg{IP}
              \item register \reg{IP} saja
              \item seluruh register umum
              \item register \reg{FLAGS} dan \reg{BP}
          \end{enumerate}

    \item Pada 8086, instruksi \asm{PUSH} bekerja dengan cara \dots
          \begin{enumerate}
              \item menambah \reg{SP} sebesar \dec{2}, lalu menyimpan
                    data
              \item mengurangi \reg{SP} sebesar \dec{2}, lalu
                    menyimpan data
              \item menyimpan data, lalu menambah \reg{SP} sebesar
                    \dec{2}
              \item menyimpan data, lalu mengurangi \reg{SP} sebesar
                    \dec{2}
          \end{enumerate}

    \item Setelah \eng{prologue} prosedur dijalankan, parameter
          pertama yang didorong paling akhir oleh pemanggil berada
          pada alamat \dots
          \begin{enumerate}
              \item \code{[BP+0]} \item \code{[BP+2]}
              \item \code{[BP+4]} \item \code{[BP+6]}
          \end{enumerate}

    \item Prosedur yang memakai \asm{RET} \code{4} berarti \dots
          \begin{enumerate}
              \item pemanggil harus menambahkan \lstasm|ADD SP, 4|
              \item prosedur membuang \dec{4} byte argumen setelah
                    mengambil alamat kembali
              \item prosedur mengembalikan nilai \dec{4} bit
              \item prosedur menyimpan empat register ke stack
          \end{enumerate}

    \item Pernyataan yang benar mengenai register \eng{callee-saved}
          pada konvensi 8086 adalah \dots
          \begin{enumerate}
              \item \reg{AX} dan \reg{DX} wajib dipulihkan prosedur
              \item \reg{BP}, \reg{SI}, dan \reg{DI} wajib dipulihkan
                    prosedur
              \item seluruh register segmen boleh dirusak prosedur
              \item \reg{SP} boleh berubah setelah \asm{RET}
          \end{enumerate}
\end{enumerate}

\subsection{Soal Uraian}

\begin{enumerate}
    \item Diketahui \reg{SS} $=$ \hex{3000} dan \reg{SP} $=$
          \hex{FFF8}. Gambarkan isi stack sel per sel (minimal lima
          sel) untuk urutan berikut: \lstasm|PUSH 000Ah|,
          \lstasm|PUSH 0014h|, \lstasm|CALL Hitung|,
          \lstasm|PUSH BP|, \lstasm|MOV BP, SP|. Tuliskan alamat
          setiap sel, isinya, dan nilai \reg{SP} serta \reg{BP}
          setelah setiap langkah.

    \item Tuliskan prosedur \code{HitungLuas} lengkap dengan
          \eng{prologue} dan \eng{epilogue}, lalu tuliskan pula
          bagian pemanggilnya. Jelaskan mengapa urutan \asm{PUSH}
          pada pemanggil menentukan nomor offset \code{[BP+4]} dan
          \code{[BP+6]}.

    \item Sebuah program berprosedur menghasilkan perilaku yang tidak
          menentu: kadang berjalan benar, kadang melompat ke alamat
          acak. Sebutkan empat kemungkinan penyebab yang berkaitan
          dengan pengelolaan stack, lalu jelaskan langkah pemeriksaan
          untuk masing-masing penyebab dengan bantuan
          \eng{debugger}.

    \item Bandingkan prosedur faktorial rekursif dan iteratif untuk
          $n = \dec{12}$ dari segi pemakaian stack, jumlah
          pemanggilan, dan risiko \eng{stack overflow}. Simpulkan
          kapan rekursi tetap merupakan pilihan yang tepat.
\end{enumerate}

\subsection{Kisi-kisi Penilaian}

\begin{table}[H]
    \centering
    \caption{Kisi-kisi evaluasi Bab~\ref{chap:prosedur}.}
    \label{tab:kisi-bab13}
    \small
    \begin{tabular}{@{}p{1.5cm}p{5.2cm}p{2.8cm}c@{}}
        \toprule
        \textbf{Sub-CPMK} & \textbf{Indikator yang diukur} &
        \textbf{Bentuk soal} & \textbf{Bobot} \\
        \midrule
        3.4 & Menjelaskan kerja \asm{CALL} dan \asm{RET} serta isi
        stack saat pemanggilan & Pilihan ganda 1, 2 & 15\% \\
        \addlinespace
        3.4 & Menganalisis \eng{stack frame} dan akses parameter
        melalui \reg{BP} & Pilihan ganda 3, 4; Uraian 2 & 30\% \\
        \addlinespace
        3.4 & Mengelola konvensi pemanggilan dan penyimpanan register
        & Pilihan ganda 5; Uraian 3 & 25\% \\
        \addlinespace
        3.4 & Menelusuri isi stack dan menilai pemakaian rekursi &
        Uraian 1, 4 & 30\% \\
        \bottomrule
    \end{tabular}
\end{table}

\begin{checklistbox}
\textbf{Checklist Pencapaian Kompetensi Bab~\ref{chap:prosedur}}

Seorang mahasiswa dinyatakan kompeten pada Sub-CPMK 3.4 apabila:
\begin{itemize}
    \item[\checked] dapat menjelaskan mekanisme \asm{CALL} dan
          \asm{RET} beserta isi stack yang terlibat;
    \item[\checked] dapat menggambar isi stack sel per sel pada
          setiap langkah pemanggilan prosedur;
    \item[\checked] dapat menulis prosedur berparameter dengan
          \eng{prologue}, \eng{epilogue}, dan \asm{RET} \code{n}
          yang benar;
    \item[\checked] dapat mengembalikan nilai prosedur melalui
          \reg{AX} serta menjelaskan pemakaian \reg{DX}:\reg{AX};
    \item[\checked] dapat menulis prosedur rekursif dengan kasus
          basis dan menjelaskan risiko \eng{stack overflow};
    \item[\checked] dapat menemukan dan memperbaiki kesalahan stack
          yang tidak seimbang pada program nyata.
\end{itemize}
\end{checklistbox}
